\section{\en{Federated Learning} και Εφαρμογές του στην Ασφάλεια Δικτύων}

\subsection{Εισαγωγή}

Η αυξανόμενη συλλογή δεδομένων στα σύγχρονα δίκτυα επικοινωνιών έχει δημιουργήσει σημαντικές ανησυχίες σχετικά με την ιδιωτικότητα, την επεκτασιμότητα και την κατανεμημένη νοημοσύνη. Στο πλαίσιο των δικτύων \en{5G} και του αναδυόμενου \en{6G}, όπου παράγονται τεράστιες ποσότητες ευαίσθητων δεδομένων από γεωγραφικά κατανεμημένες συσκευές, οι παραδοσιακές κεντρικοποιημένες προσεγγίσεις μηχανικής μάθησης παρουσιάζουν σημαντικούς περιορισμούς. Το \en{Federated Learning} έχει αναδειχθεί ως ένα υποσχόμενο εργαλείο που επιτρέπει τη συνεργατική εκπαίδευση μοντέλων χωρίς ανταλλαγή δεδομένων. Διατηρώντας τα δεδομένα τοπικά και κοινοποιώντας μόνο ενημερώσεις μοντέλων, το \en{FL} αντιμετωπίζει ζητήματα ιδιωτικότητας και επιτρέπει κατανεμημένη μάθηση μεγάλης κλίμακας \cite{kairouz2021advances}.

\subsection{Ορισμός και Κίνητρο του \en{Federated Learning}}

Το \en{Federated Learning} είναι ένα κατανεμημένο πλαίσιο μηχανικής μάθησης όπου πολλοί πελάτες (\en{clients}) εκπαιδεύουν συλλογικά ένα κοινό μοντέλο υπό τον συντονισμό ενός κεντρικού εξυπηρετητή (\en{server}), διατηρώντας τα δεδομένα τους ιδιωτικά \cite{mcmahan2017communication}. Τα βασικά κίνητρα περιλαμβάνουν την προστασία της ιδιωτικότητας των δεδομένων, καθώς αυτά δεν μπορούν να κοινοποιηθούν εύκολα, την κατανεμημένη φύση των δεδομένων που βρίσκονται σε πολλές συσκευές, καθώς και τους περιορισμούς επικοινωνίας, αφού η μεταφορά μεγάλων όγκων δεδομένων είναι ιδιαίτερα δαπανηρή.

\subsection{Κεντρικοποιημένη έναντι Κατανεμημένης Μάθησης}

Η παραδοσιακή μηχανική μάθηση είναι κεντρικοποιημένη, συγκεντρώνοντας δεδομένα σε ένα σημείο, γεγονός που δημιουργεί κινδύνους ασφάλειας και ιδιωτικότητας. Αντίθετα, το \en{FL} είναι αποκεντρωμένο, διατηρώντας τα δεδομένα τοπικά και κοινοποιώντας μόνο ενημερώσεις μοντέλων. Με τον τρόπο αυτό μειώνεται ο κίνδυνος διαρροής δεδομένων, διευκολύνεται η συμμόρφωση με κανονισμούς όπως ο \en{GDPR} και επιτυγχάνεται καλύτερη επεκτασιμότητα. Ωστόσο, εισάγονται νέες προκλήσεις, όπως η ανομοιογένεια των δεδομένων και τα προβλήματα επικοινωνίας.

\subsection{Διαδικασία \en{Federated Learning}}

Η διαδικασία του \en{Federated Learning} περιλαμβάνει επαναλαμβανόμενους κύκλους συνεργασίας μεταξύ εξυπηρετητή και πελατών. Αρχικά, ο εξυπηρετητής διανέμει το αρχικό μοντέλο στους πελάτες. Στη συνέχεια, κάθε πελάτης πραγματοποιεί τοπική εκπαίδευση χρησιμοποιώντας τα δικά του δεδομένα. Έπειτα, οι ενημερώσεις αποστέλλονται πίσω στον εξυπηρετητή, ο οποίος τις συνδυάζει μέσω τεχνικών όπως το \en{Federated Averaging} \cite{mcmahan2017communication}. Η διαδικασία επαναλαμβάνεται μέχρι να επιτευχθεί σύγκλιση του μοντέλου.

\subsection{Ιδιωτικότητα και Ασφάλεια στο \en{FL}}

Το \en{Federated Learning} προσφέρει σημαντικά πλεονεκτήματα ιδιωτικότητας, καθώς τα δεδομένα δεν εγκαταλείπουν ποτέ τις τοπικές συσκευές, μειώνοντας έτσι τον κίνδυνο διαρροής. Παρ’ όλα αυτά, παραμένουν σημαντικοί κίνδυνοι, όπως οι επιθέσεις δηλητηρίασης μοντέλου, όπου κακόβουλοι πελάτες επιχειρούν να αλλοιώσουν τη συνολική συμπεριφορά του μοντέλου \cite{bhagoji2019analyzing}, καθώς και η διαρροή πληροφοριών μέσω των \en{gradients}, από τα οποία μπορεί να ανακτηθούν ευαίσθητα δεδομένα \cite{zhu2019deep}. Για την αντιμετώπιση αυτών των απειλών χρησιμοποιούνται τεχνικές όπως η Διαφορική Ιδιωτικότητα (\en{Differential Privacy}) \cite{dwork2014algorithmic} και η Ασφαλής Συνάθροιση (\en{Secure Aggregation}) \cite{bonawitz2017practical}.

\subsection{Εφαρμογές στην Ασφάλεια Δικτύων}

Το \en{Federated Learning} βρίσκει σημαντικές εφαρμογές στην ασφάλεια δικτύων, ιδιαίτερα σε συστήματα ανίχνευσης εισβολών (\en{Intrusion Detection Systems -- IDS}) και στην κατανεμημένη ανίχνευση ανωμαλιών. Στα δίκτυα τηλεπικοινωνιών προσφέρει αυξημένη επεκτασιμότητα, επιτρέπει συνεργασία μεταξύ διαφορετικών φορέων χωρίς κοινοποίηση ευαίσθητων δεδομένων και ενισχύει τη νοημοσύνη στην περιφέρεια του δικτύου (\en{Edge Intelligence}), επιτρέποντας την άμεση επεξεργασία δεδομένων και τη λήψη αποφάσεων σε τοπικό επίπεδο.

\subsection{Πλαίσιο \en{Flower} για \en{Federated Learning}}

Το \en{Flower} είναι ένα ανοικτού κώδικα πλαίσιο για την ανάπτυξη συστημάτων \en{Federated Learning} \cite{beutel2020flower}. Η αρχιτεκτονική του βασίζεται σε έναν εξυπηρετητή (\en{server}) που συντονίζει τη διαδικασία και σε πολλούς πελάτες (\en{clients}) που εκτελούν τοπική εκπαίδευση. Μεταξύ των βασικών χαρακτηριστικών του περιλαμβάνονται η ευελιξία, η συμβατότητα με πλαίσια όπως \en{PyTorch} και \en{TensorFlow}, καθώς και η υψηλή επεκτασιμότητα. Επιπλέον, το \en{Flower} θεωρείται ιδιαίτερα κατάλληλο για την κατανεμημένη εκπαίδευση \en{LLMs}, ειδικά όταν συνδυάζεται με αποδοτικές τεχνικές όπως το \en{LoRA}.
